% Generated draft: requires mathematical and source review.
\documentclass[11pt]{article}
\usepackage[margin=25mm]{geometry}
\usepackage{fontspec}
\setmainfont{Latin Modern Roman}
\usepackage{amsmath,amssymb,mathtools,unicode-math}
\setmathfont{Latin Modern Math}
\usepackage{enumitem,needspace,xurl,hyperref}
\setlength{\parindent}{0pt}
\setlength{\parskip}{5pt}
\setlength{\emergencystretch}{3em}
\newcommand{\N}{\mathbb{N}}
\newcommand{\Z}{\mathbb{Z}}
\newcommand{\Q}{\mathbb{Q}}
\newcommand{\R}{\mathbb{R}}
\newcommand{\C}{\mathbb{C}}
\newcommand{\F}{\mathbb{F}}
\newcommand{\A}{\mathbb{A}}
\newcommand{\PP}{\mathbb{P}}
\newcommand{\E}{\mathbb{E}}
\newcommand{\eps}{\varepsilon}
\newcommand{\dd}{\,\mathrm{d}}
\begin{document}
\section*{UnsolvedMath ID: 30002816}
\textbf{Source problem code:} catalogue entry
\subsection{Definitions and mathematical statement}
The supplied record asks whether every calibration of every Grassmann orbitope admits a sum-of-squares certificate with linear summands. It does not specify the admissible integer pairs \((k,n)\), the ambient coordinates, or the precise vanishing ideal. Accordingly, the following is a conditional formulation of the requested certificate, pending those specifications.

Under a real-coordinate interpretation, let \(V_{k,n}\) be the intended finite-dimensional real ambient vector space, let \(G_{k,n}\subseteq V_{k,n}\) be the intended Grassmann orbitope, and let \(\R[V_{k,n}]\) denote its ambient polynomial ring. Let \(I_{k,n}\subseteq\R[V_{k,n}]\) be the intended vanishing ideal. For a linear functional \(\ell\in V_{k,n}^{*}\), a supporting value \(\delta\in\R\) satisfies
\[
\ell(p)\leq\delta\quad\text{for every }p\in G_{k,n},
\qquad
\ell(p_0)=\delta\quad\text{for some }p_0\in G_{k,n}.
\]
Thus \(f=\delta-\ell\) is nonnegative on \(G_{k,n}\) and vanishes at a supporting point.

\textbf{Certificate question.} For every admissible \((k,n)\) and every supporting pair \((\delta,\ell)\) corresponding to a calibration in the source's convention, do there exist a finite integer \(m\geq0\) and permitted linear polynomials \(h_1,\ldots,h_m\) such that
\[
\delta-\ell-\sum_{i=1}^{m}h_i^2\in I_{k,n}?
\]
Equivalently, writing brackets for classes in the quotient ring, the requested identity is
\[
[\delta-\ell]=\sum_{i=1}^{m}[h_i]^2
\quad\text{in }\R[V_{k,n}]/I_{k,n}.
\]
If ``linear'' means affine linear, the permitted representatives are \(h_i(p)=a_i+\lambda_i(p)\), with \(a_i\in\R\) and \(\lambda_i\in V_{k,n}^{*}\). If it means homogeneous linear, one must instead require \(a_i=0\). The supplied record does not resolve this distinction.

\textbf{Unresolved qualifications.} The source also writes the certificate as a pointwise identity \(\delta-\ell(p)=\sum_i h_i(p)^2\) on \(G_{k,n}\). A congruence modulo a vanishing ideal gives an identity on the set defining that ideal; its identification with an identity on the entire orbitope requires clarification. In particular, the ideal of the convex orbitope and the ideal of a Grassmannian or generating orbit must not be identified without justification. The source does not define ``calibration,'' specify its normalization, or establish here which supporting pairs it includes. Its assertion that this question is equivalent to Harvey--Lawson Question 6.5 is retained as a supplied attribution, not as an independently verified equivalence.
\subsection{Short English statement}
For every Grassmann orbitope, the supplied question asks whether each supporting nonnegative function associated with a calibration is a sum of squares of linear polynomial classes modulo the intended vanishing ideal. A precise unconditional statement requires the parameter range, coordinate model, calibration convention, meaning of ``linear,'' and defining set of the ideal. The record labels the problem open; this draft preserves that supplied label without independently verifying it.
\subsection{Sources}
\begin{itemize}
\item Supplied local UnsolvedMath registry record, permanent ID \(30002816\), title \emph{Linear Sum-of-Squares Representations on Grassmann Orbitopes}, exact source problem code \texttt{catalogue entry}. Project registry locations: \texttt{data/unsolvedmath/problems.json} and supplementary \texttt{catalog/problems.json}. The supplied record belongs to \emph{Oberwolfach Reports Open Problems} and gives proposed year 2015; no individual proposer is supplied.
\item \emph{Discrete Differential Geometry} (2015), Oberwolfach Reports, DOI 10.4171/owr/2015/13, pp.\ 661--729. Supplied source report and literature-source link: \url{https://doi.org/10.4171/owr/2015/13}. The record describes these reports as documenting workshops and their open-problem discussions. No precise page or report contribution locating this question is supplied.
\item Sanyal--Sottile--Sturmfels, supplied literature-source link: \url{https://arxiv.org/abs/0911.5436}. The record does not supply a title or a passage establishing the relevant formulation.
\item Literature assessment dated exactly \texttt{2026-08-21}: the supplied curation labels the problem open and states that the all-Grassmann-orbitope conjecture, described there as equivalent in the report to Harvey--Lawson Question 6.5, has no complete proof or disproof in the checked primary literature. This is supplied dated curation, not independently verified literature assessment. The supplied background expressly requests independent verification of time-sensitive status and cited claims.
\item The registry labels the statement \texttt{corrected\_verified} and reports manual semantic, typographic, and LaTeX review of a meaning-preserving cleanup. That supplied editorial verification is not independent mathematical verification of this draft.
\end{itemize}
\end{document}
